\documentclass[11pt]{article}

% Standalone manuscript. The earlier draft depended on an ICLR style file that
% is not present in this repository; this version compiles as a standard article.
\usepackage[margin=1in]{geometry}
\usepackage{lmodern}
\usepackage{amsmath,amssymb}
\usepackage{booktabs}
\usepackage{graphicx}
\usepackage{array}
\usepackage{hyperref}
\hypersetup{colorlinks=true,linkcolor=blue,citecolor=blue,urlcolor=blue}
\emergencystretch=2em

\newcommand{\care}{\textsc{CARE}}
\newcommand{\botans}{\bot}
\newcommand{\gold}{\mathcal{G}}
\newcommand{\sources}{\mathcal{S}}

\title{\care: Conditional Answer Repair with Source-Anchored Evidence\\
for Long-Horizon Conversational Memory}
\author{Anonymous authors}
\date{}

\begin{document}
\maketitle

\begin{abstract}
Long-horizon conversational question answering is often framed as a problem
of learning what an agent should store. Yet a more immediate failure arises
when the original conversation is available in external storage but the
answering model receives the wrong evidence, loses speaker or temporal
attribution, or sees only a truncated context. We present \care{} (Conditional
Answer Repair with Evidence), a source-anchored read path that preserves
immutable, addressable conversation turns. It combines inexpensive sparse
retrieval with a reader that returns an answer, supporting source identifiers,
and targeted searches for missing evidence. Deterministic checks on source
membership and selected person, event, and temporal constraints trigger a
single broader retrieval-and-answer pass only when needed. On all 1,986
questions from ten LoCoMo conversations, using the same pinned Qwen3-8B
answer model, \care{} obtains a LoCoMo aggregate score of 0.5256, compared
with 0.4069 for dense similarity retrieval over original turns, 0.1961 for
an 8K-window context baseline, and 0.0774 without memory. Its raw
BERTScore F1 is 0.9088. The paired aggregate gain over dense retrieval is
0.1187 (conversation-cluster bootstrap 95\% interval: [0.0837, 0.1503]);
\care{} also raises annotated-evidence recall from 0.5753 to 0.7033.
A post-hoc coverage analysis qualifies the context result: the 8K baseline
retains only 25.8\% of turns and no annotated answer evidence for 1,060
answerable questions, but outperforms \care{} when all annotated evidence
remains visible. These findings identify conditional, source-aware reading
as a useful component of agent memory while delimiting the claim: this
experiment does not establish an advantage over untruncated full context
or demonstrate learned self-evolution of stored memories.
\end{abstract}

\section{Introduction}
\label{sec:introduction}

An agent that answers questions about months of conversation faces a
selection problem. Its historical record can be longer than the model's
working context; a fact relevant to the current question may have occurred
in an early session, whereas a recent exchange may be lexically similar but
refer to a different person or event. LoCoMo was designed to expose such
long-range conversational failures across sessions, temporal references,
and questions with unsupported premises~\cite{maharana2024locomo}. The
usual responses are to extend the context window or to write a compact
external memory. Both still require the answer model to see, identify, and
use the correct evidence at the moment of answering.

Research on long-term agent memory has proposed atomic notes, learned write
operations, dynamic links, value-aware retrieval, and runtime utility
updates. These mechanisms concern how a memory \emph{store} should evolve.
The experiments in this paper isolate a complementary prerequisite:
\emph{How can an agent recover a supportable answer from an immutable
multi-session record under a bounded reading context?} A learned memory
writer can help only if the downstream reader preserves the identity of the
person, event, time, and source that make an answer valid. Conversely,
compressing or rewriting a source may erase details needed for a later
question. We therefore study the read path over original turns and reserve
claims about self-evolving stores for future experiments.

Our answer is \care{}, a conditional read-and-repair procedure. Each turn is
an addressable source with its session date and speaker. A first reader call
uses a small sparsely retrieved bundle and emits both a concise answer and
source identifiers. If the answer lacks an admissible citation, contradicts
simple person or event checks, or explicitly identifies an evidence gap,
\care{} expands the candidate set using lexical, stemmed, and local-neighbor
views, then asks the reader once more. This design makes provenance part of
the inference contract. It avoids paying for broad evidence on every
question, while giving the reader a bounded opportunity to repair a narrow
first pass.

The empirical result is also conditional. On a complete, same-backbone
1,986-question comparison, \care{} improves the LoCoMo aggregate score over
dense turn retrieval by 0.1187. It makes 2,990 answer-model calls for 1,986
questions, including 1,004 repair calls. However, an 8K context baseline
performs better than \care{} on the subset whose annotated answer sources
fit inside its retained context. Its poor overall score is strongly
associated with losing earlier evidence, not evidence that direct reading
is intrinsically inferior. This distinction matters when positioning a
memory method against long-context models.

\paragraph{Contributions.}
\begin{enumerate}
  \item We formulate and implement a \emph{source-anchored conditional read
  path}: immutable turns, explicit citation identifiers, bounded first-pass
  retrieval, deterministic support checks, and one risk-gated repair pass.
  The contribution is an answer-time memory interface, rather than an
  unmeasured learned writer or a new backbone model.
  \item We provide a full matched-backbone evaluation against question-only,
  8K context, and dense raw-turn retrieval baselines on all 1,986 LoCoMo
  questions, reporting the task-specific score, BERTScore, annotated-evidence
  recall, per-category behavior, and model-call counts.
  \item We conduct a source-coverage-conditioned context analysis. It shows
  where an 8K context baseline fails because annotated evidence is outside
  the window and where direct reading succeeds when evidence is visible.
  \item We report a negative relation-and-temporal extension result and
  distinguish complete comparisons from earlier partial graph and BM25
  cohorts. This narrows the experimentally supported claim and identifies
  the next controls needed to study self-evolving memory.
\end{enumerate}

\section{Related Work and Research Position}
\label{sec:related}

Retrieval-augmented generation accesses external text at answer time rather
than relying solely on model parameters~\cite{lewis2020rag}. The quality of
an answer can therefore depend separately on retrieval coverage and on the
reader's ability to use retrieved passages. For multi-session agents,
MemGPT uses a hierarchy of context and external memory to manage a limited
working window~\cite{packer2023memgpt}. Mem0 explores extraction,
consolidation, retrieval, and a graph-enhanced variant for long-term agent
memory~\cite{chhikara2025mem0}. A-Mem creates structured, dynamically linked
notes and revises contextual attributes as new information arrives
~\cite{amem2025}. Memory-R1 learns write operations and an answering policy
with reinforcement learning~\cite{memoryr12025}. These systems motivate the
broader ASEM draft's concern with persistent memory organization.

\care{} investigates a complementary design point. It leaves original turns
immutable and places adaptation in the \emph{read operation}: source-aware
retrieval, an answer that carries admissible citations, and conditional
evidence repair. This is distinct from claiming that a graph, write policy,
or learned utility estimate produced the observed gain. It is also distinct
from a conventional dense top-$k$ reader, which retrieves semantically close
turns but does not ask for missing-source searches or validate person and
event bindings. The distinction is operational rather than a claim that no
prior work has considered evidence-grounded answers.

LoCoMo provides a useful setting because its annotated questions refer to
long, multi-session conversations and include unanswerable questions
~\cite{maharana2024locomo}. We use its answer annotations and source IDs
only for post-hoc evaluation. BERTScore~\cite{zhang2020bertscore} complements
the task metric by measuring textual semantic overlap; it cannot determine
whether an answer is supported by a conversation.

\section{Problem Setting}
\label{sec:setting}

Let a conversation be a sequence of original turns
\begin{equation}
\sources = (s_1,\ldots,s_N),\qquad
s_i=(\mathrm{id}_i,\mathrm{session}_i,t_i,p_i,x_i),
\label{eq:source}
\end{equation}
where $t_i$ is the session date, $p_i$ the speaker, and $x_i$ the turn text
with an available image caption appended as text. The source identifier
$\mathrm{id}_i$ is unique within a conversation. A question $q$ asks for a
short answer $a$ or the abstention symbol $\botans$ if the premise lacks
support. The reader has a bounded prompt window $W$, so a source-selection
procedure must produce a visible bundle $B(q)\subseteq\sources$ before
generation. The objective is to maximize answer quality while retaining
source provenance and limiting model calls.

This setting does not learn a write policy, update note embeddings, or train
a utility function.
All methods in the main comparison receive the same raw conversation history
from which to construct their answer context; their mechanisms for selecting
that context differ.

\section{Method: Conditional Source-Anchored Reading}
\label{sec:method}

\subsection{Immutable Source Index and Cheap First Pass}

For each conversation, \care{} indexes every turn by its source ID, session,
speaker, and text. A BM25-style sparse score ranks full turns against the
question, and the initial bundle contains the eight highest-ranked turns.
The ranking is deterministic, with source order breaking ties. The retained
source record preserves who spoke and when, which is essential for questions
containing first-person pronouns or relative dates. A first answer call sees
the question and a compact serialization of these original turns, rather
than an LLM-generated summary whose source binding may be lost.

\subsection{An Answer Contract That Exposes Gaps}

The reader returns a JSON object
\begin{equation}
(a,C,Q)=R_{\theta}(q,B),
\label{eq:reader}
\end{equation}
where $a$ is the proposed short answer, $C$ is a list of supporting source
IDs, and $Q$ contains at most two short searches for missing evidence. A
supported answer should cite visible turns; an unsupported premise should
produce ``No information available.'' The prompt asks the reader to check
the named person, event, and time before answering and to resolve relative
dates only when an explicit session date supports the calculation.

The returned citation is an \emph{audit handle}, not a proof of entailment.
\care{} verifies that every cited ID is visible and applies deterministic
checks to guard against assigning a speaker's first-person claim to another
named person, confusing explicit event anchors, or citing a turn that lacks
a requested cessation event. These checks reject some readily identifiable
failures; they cannot certify all semantic relations. If a generated object
is malformed or a supported answer has no admissible citation, the answer
is converted to abstention for that pass.

\subsection{Risk-Gated Evidence Expansion}

Let $V(q,a,C,B)$ be the result of the parse, visibility, and source-binding
checks above. A second pass is triggered by
\begin{equation}
\tau(q)=\mathbf{1}\!\left[
\neg V(q,a,C,B_1)\ \lor\
\bigl(a=\botans\ \land\ (Q\neq\varnothing\ \lor\ \rho(q,B_1))\bigr)
\right],
\label{eq:gate}
\end{equation}
where $\rho$ is a simple risk flag for temporal or multi-person wording and
update cues. When $\tau=1$, \care{} assembles a broader source pool from
the top 18 lexical turns, stemmed sparse ranking, neighbors of leading
turns within the same session, and searches proposed by the first reader.
The expanded reader is given a correction to check direct evidence for the
person, event, and time in the question. Both passes use original turns;
the visible evidence serialization is capped at 10,000 characters. The
procedure makes at most two answer-model calls per question.

\subsection{Relation-and-Temporal Extension}

Explicit links and temporal consistency are recurring goals in long-term
memory research. We therefore evaluate a conservative extension that adds relation and
temporal proof checks to the same read path. It is an empirical ablation,
not a component of the main \care{} result. The extension changes only 15
of 1,986 answers and does not improve the aggregate score. Reporting this
negative result avoids attributing gains to additional structure that the
complete experiment does not support.

\section{Experimental Design}
\label{sec:experiments}

\subsection{Dataset and Fixed Inference Configuration}

We evaluate question answering on the ten LoCoMo conversations available in
the locked dataset~\cite{maharana2024locomo}: 5,882 conversation turns and
1,986 questions. The question categories contain 282 single-hop, 321
temporal, 96 commonsense/inference, 841 conversational-context, and 446
adversarial or unsupported-premise questions. The adapter passes text,
session dates, speaker labels, and existing image captions. The experiment
does not process image pixels and should not be presented as a multimodal
benchmark result.

All primary arms use Qwen3-8B~\cite{yang2025qwen3} at the pinned repository
revision \texttt{b968826d9c46dd60\allowbreak 66d109eabc625518\allowbreak 8de91218}, served with
vLLM on two NVIDIA T4 GPUs. The context window is 8,192 tokens; decoding
uses temperature zero, a maximum of 224 answer tokens, and Qwen's
non-thinking mode. Five shards contain whole conversations. The same 1,986
question IDs and scorer are verified before merging each arm, and all arms
finish without inference errors.

\subsection{Baselines}

\paragraph{NoMemory.} The backbone sees the question without the historical
conversation. Its earlier baseline prompt requests a best guess when
uncertain; this affects unsupported-premise questions.

\paragraph{FullContext-8K.} The reader receives dated original turns. For
each conversation, a fixed, query-independent context keeps the first five
turns and the longest recent suffix that fits the 8K budget after reserving
answer tokens. This is a \emph{window-limited} context baseline, not a model
that sees the entire multi-session transcript.

\paragraph{SimRetrieval.} A pinned
\texttt{sentence-transformers/all-MiniLM-L6-v2} encoder embeds original
dated turns and the question. Mean-pooled, $\ell_2$-normalized vectors
retrieve the top 18 turns by cosine similarity, with a 10,000-character
evidence cap. The answer model is called once. This is a raw-turn dense
retrieval control; it does not reproduce the older project's note-writing
SimRetrieval implementation.

\paragraph{Historical graph and BM25 controls.} Earlier Qwen3-8B graph,
source-hybrid, and BM25 runs are used only on their available, exactly
paired question IDs. They are not part of the full four-arm table because
the old graph run completed only 746 overlapping questions and the BM25
checkpoint only 1,128; prompts and budgets also differ.

\subsection{Metrics and Statistical Comparison}

We implement the LoCoMo answer scorer: category 1 averages the best
stemmed-token F1 for each comma-separated reference item; categories 2--4
use stemmed-token F1, with the category-3 reference truncated before a
semicolon; category 5 assigns one if the response contains one of the
benchmark's accepted refusal phrases (``no information available'' or
``not mentioned'') and zero otherwise. The \emph{aggregate LoCoMo score}
averages these category-specific values; it is not token F1 on all items.
We report answerable and category-5 groups separately.

We also report raw BERTScore F1 and its baseline-rescaled variant. Both use
the same \texttt{roberta-large} revision, layer 17, no IDF, and identical
reference policy across arms~\cite{zhang2020bertscore}. The main tables
show raw F1. Annotated-evidence recall is the fraction of gold source IDs
present in the visible answer bundle, averaged over answerable questions.
For paired LoCoMo differences we resample the ten conversations, rather
than individual questions, 10,000 times and report percentile 95\%
intervals. These intervals describe this public, small-cluster cohort; they
are not evidence of performance on an independent held-out population.

\section{Results}
\label{sec:results}

\subsection{Full Same-Backbone Comparison}

\begin{table}[t]
\centering
\caption{Complete LoCoMo comparison on 1,986 identical questions. The
aggregate mixes token-based scores for categories 1--4 with binary
refusal accuracy for category 5. BERTScore is raw F1.}
\label{tab:main}
\small
\begin{tabular}{lrrrr}
\toprule
Method & Aggregate & Answerable F1 & Refusal correct & BERTScore F1 \\
\midrule
NoMemory & 0.0774 & 0.0998 & 0/446 & 0.8776 \\
FullContext-8K & 0.1961 & 0.2406 & 19/446 & 0.8912 \\
SimRetrieval (raw turns) & 0.4069 & 0.3124 & 327/446 & 0.8896 \\
\care{} Bridge v2 & \textbf{0.5256} & \textbf{0.4648} & 328/446 & \textbf{0.9088} \\
Relation+Temporal extension & 0.5251 & 0.4603 & \textbf{334/446} & 0.9083 \\
\bottomrule
\end{tabular}
\end{table}

Table~\ref{tab:main} shows a 0.1187 aggregate advantage for \care{} over
raw-turn dense retrieval, with a conversation-cluster 95\% interval of
[0.0837, 0.1503]. The answerable difference is 0.1524
([0.1287, 0.1770]). The difference on category-5 refusal is one question
(328 versus 327), with a cluster interval for the accuracy difference that
includes zero. Thus the main observed benefit over dense retrieval is in
answering supported questions. \care{} reaches 0.7033 annotated-evidence
recall versus 0.5753 for dense retrieval. Its conditional repair requires
2,990 model calls for 1,986 questions, 50.6\% more than the single-call
baselines. The relation-and-temporal extension makes 3,000 calls and
changes 15 answers, lowering the aggregate from 0.5256 to 0.5251.

The NoMemory and FullContext-8K refusal scores need particular care:
NoMemory was prompted to guess, whereas \care{} and SimRetrieval were
instructed to abstain when evidence was missing. The category-5 gap
therefore partly reflects the answer contract. BERTScore also must not
substitute for refusal accuracy: it measures semantic overlap with a
standard ``Not mentioned'' reference even if the response is not a valid
benchmark refusal.

\subsection{Performance by Question Type}

\begin{table}[t]
\centering
\caption{Category-specific LoCoMo scores. Category 5 is binary refusal
accuracy; other cells follow the benchmark's token-based scorer.}
\label{tab:category}
\small
\begin{tabular}{lrrrrr}
\toprule
Category & $n$ & NoMemory & FullContext-8K & SimRetrieval & \care{} \\
\midrule
Single-hop (1) & 282 & 0.0966 & 0.2353 & 0.2283 & \textbf{0.2609} \\
Temporal (2) & 321 & 0.0336 & 0.2367 & 0.3019 & \textbf{0.4415} \\
Commonsense (3) & 96 & 0.1256 & \textbf{0.2014} & 0.1263 & 0.1647 \\
Conversational (4) & 841 & 0.1232 & 0.2484 & 0.3658 & \textbf{0.5763} \\
Unsupported (5) & 446 & 0.0000 & 0.0426 & 0.7332 & \textbf{0.7354} \\
\bottomrule
\end{tabular}
\end{table}

The largest \care{} gains over dense retrieval occur on temporal and
conversational-context questions. This pattern is consistent with a benefit
from finding and binding old source turns, although the read prompt and
validation rules differ and the experiment does not isolate either factor.
Category 3 remains a weakness: FullContext-8K scores 0.2014 versus 0.1647
for \care{}. In the older paired graph cohort, the previous graph reader
also does better than \care{} on category-1 list questions and on refusal
accuracy. These exceptions rule out a claim of uniform superiority.

\subsection{The Context-Coverage Crossover}

The 8K baseline retains 1,515 of 5,882 turns across the ten conversations
(25.8\%); all conversations are truncated, with per-conversation retention
between 21.9\% and 43.9\%. Its mean server-reported prompt length is
approximately 7,879 tokens, close to the configured window. All 1,986
completions stop normally, so its low score is not a generation-length
failure. Dense retrieval uses approximately 1,112 prompt tokens per
question on average, and NoMemory approximately 68. We do not report a
comparable total prompt-token cost for \care{} because its per-call token
receipts were not retained in the merged artifact.

To distinguish evidence availability from answer quality, we reconstruct
the fixed FullContext-8K selection and compare it to LoCoMo's annotated
source IDs \emph{after} inference. Among 1,540 answerable questions, nine
have malformed or absent source IDs and four lack annotations; we exclude
these 13 from the coverage-conditioned table. The resulting 1,527
questions are partitioned into all, partial, or no annotated sources
retained by the 8K context.

\begin{table}[t]
\centering
\caption{Answerable score conditioned on post-hoc annotated-source
coverage in the fixed 8K context. Source annotations are diagnostic and
were never available to a model during inference.}
\label{tab:coverage}
\small
\begin{tabular}{lrrrr}
\toprule
Annotated sources in 8K & $n$ & FullContext-8K & SimRetrieval & \care{} \\
\midrule
All retained & 337 & \textbf{0.5926} & 0.3344 & 0.4939 \\
Partly retained & 130 & \textbf{0.3242} & 0.1941 & 0.2152 \\
None retained & 1,060 & 0.1191 & 0.3211 & \textbf{0.4892} \\
\bottomrule
\end{tabular}
\end{table}

When all annotated sources remain visible, FullContext-8K exceeds \care{}
by 0.0987 F1; the paired cluster interval for \care{} minus context is
[-0.1378, -0.0620]. When no annotated source remains, \care{} exceeds the
context baseline by 0.3701 ([0.3309, 0.4079]). This crossover is a central
result: source-directed memory access is valuable when a bounded context
omits answer evidence, while direct reading can be stronger when the
evidence fits. An annotation may be incomplete; ``no annotated source''
does not prove that every alternative supporting phrase is absent.

\subsection{Earlier Incomplete Comparisons}

Earlier Qwen3-8B graph and BM25 predictions were rescored and paired by
exact question ID, but their runs did not cover the full set. On 746
overlapping questions, \care{} scores 0.5188 versus 0.3786 for the old graph
reader and 0.3467 for the source-hybrid reader. On 1,128 overlapping
questions it scores 0.4912 versus 0.1372 for the old BM25 reader.
These are cohort-specific, whole-pipeline comparisons; differing prompts,
budgets, and refusal wording prevent an isolated architecture claim.
The previous graph and source-hybrid readers correctly refuse more
category-5 questions than \care{} within their 156-question overlap.

An earlier full 1,986-question FullContext result used Qwen3-4B and a
different scoring protocol; its reported 0.3485 F1 should not be placed in
Table~\ref{tab:main}. A separate earlier Qwen3-8B FullContext attempt
exceeded a 4,096-token evidence allowance and did not produce usable
answers. Neither is evidence about the quality of an untruncated
Qwen3-8B FullContext reader under the present evaluation contract.

\section{Discussion}
\label{sec:discussion}

\paragraph{What the experiments support.}
The evidence supports a practical claim about the \emph{read path}:
explicit source handles, speaker and date preservation, and a conditional
chance to search again can improve answer quality under an 8K budget.
The gain over dense retrieval is paired on every question and accompanied
by higher annotated-source recall. It is obtained without training or
changing Qwen3-8B parameters. The context-coverage crossover sharpens this
claim: the method's comparative advantage is strongest when the relevant
source is outside a fixed context window.

\paragraph{What remains unproven.}
The evaluation does not establish that \care{} is better than a model that
reads all 5,882 turns, nor does it identify how much of the gain comes from
sparse retrieval, the reader's citation contract, the abstention phrase,
or the second call in isolation. The dense control and \care{} have
different answer prompts. The NoMemory prompt's best-guess instruction
especially confounds unsupported-question comparisons. Category-5 scoring
checks specific phrases, and BERTScore is not a factual verifier. The ten
public conversations were also used during method development, so these
scores are not an independent generalization test.

\paragraph{Relation to evolving memory stores.}
Multi-attribute notes, learned write decisions, bidirectional links, and
value-aware retrieval remain plausible ways to
improve a long-term store. The present full experiment does not exercise
those mechanisms. It instead identifies a source-anchored reader that could
serve as their evaluation-preserving front end: a future writer can propose
compressed memories while the original turn IDs remain available for
verification and repair. Claims that memory \emph{self-evolution} caused
the reported improvement would require a controlled write-path ablation
on fresh conversations.

\paragraph{Next decisive experiments.}
The strongest context control would run the same Qwen3-8B reader with a
window large enough to retain every turn and an aligned refusal contract.
A shared-reader ablation should then vary only the candidate selector
(dense, sparse, and conditional repair), holding prompt and evidence
budget fixed. Finally, fresh histories, real image inputs, category-3
reasoning, list aggregation, and prompt-token accounting are needed for
a broader quality and efficiency claim.

\section{Conclusion}
\label{sec:conclusion}

\care{} turns long-term conversation access into a bounded, source-aware
answering procedure. On a complete 1,986-question LoCoMo comparison with
the same Qwen3-8B backbone, it improves the aggregate task score from
0.4069 for dense raw-turn retrieval to 0.5256, while using a second
answer-model call only when an initial answer or its evidence is at risk.
The accompanying context audit shows why such a reader helps: an 8K fixed
history loses most annotated answer evidence, whereas direct context
reading can be better whenever that evidence remains visible. This is a
measured contribution to the read stage of agent memory and a concrete
foundation for future, separately evaluated self-evolving stores.

\appendix
\section{Reproducibility and Score Provenance}
\label{app:reproducibility}

The locked LoCoMo JSON used in this study has SHA-256 checksum
\texttt{79fa87e90f040813\allowbreak 43b8c8debecb80a9\allowbreak a6842b76a7aa537d\allowbreak c9fdf651ea698ff4}.
The full \care{}, context, and similarity-retrieval merged reports are
stored at \path{build/care/relation_full_merged.json},
\path{build/care/context_baselines_merged.json}, and
\path{build/care/sim_retrieval_merged.json}. Post-hoc BERTScore results
are stored at corresponding \path{*_bertscore.json} paths. Exact question
ID and answer coverage, score recomputation, model revisions, and
BERTScore input hashes are validated by the merge and summary scripts.
The analysis of annotated-source visibility and cluster intervals is
reproducible with \path{scripts/summarize_care_overview.py}; its
machine-readable output is
\path{analysis/ASEM_CARE_overview_20261112_data.json}.

The answer model was pinned to \texttt{Qwen/Qwen3-8B} revision
\texttt{b968826d9c46dd60\allowbreak 66d109eabc625518\allowbreak 8de91218}.
SimRetrieval used \texttt{all-MiniLM-L6-v2} revision
\texttt{1110a243fdf4706b\allowbreak 3f48f1d95db1a4f5\allowbreak 529b4d41}.
BERTScore used \texttt{roberta-large} revision
\texttt{722cf37b1afa9454\allowbreak edce342e7895e588\allowbreak b6ff1d59}, layer 17,
without IDF. The full BERTScore reference for category 5 was
``Not mentioned''; category-3 references were truncated before the first
semicolon. Nine answerable questions with invalid source annotations
and four without source annotations were excluded only from the
coverage-conditioned analysis, not from the benchmark score.

\bibliographystyle{plain}
\bibliography{asem}
\end{document}
